\section{Discussion and Limitations}
\label{sec:discussion}

The demonstration sessions described in Section~\ref{sec:results}
support several broader observations.

\subsection{Observations}

The integrated prototype confirms that a single low-cost
microcontroller node, when paired with a small Python bridge and a
cross-platform mobile application, can deliver continuous visibility
into the four sensing modalities most often addressed by
smart-refrigerator systems. The architectural separation of the
sensor node from the camera node simplifies firmware on either side
and allows individual subsystems to be replaced without disturbing
the rest. The observed end-to-end latency, measured from the
Arduino's serial output to the corresponding update in the Flutter
Dashboard, was below one second; this margin is more than adequate
for inventory-monitoring use cases in which the underlying physical
processes evolve over minutes to days.

\subsection{Comparison with Related Work}

Relative to camera-only smart-refrigerator
systems~\cite{ericyang2022smartfridge,vadlamudi2025fridgesense}, the
proposed prototype does not depend on a trained classifier and
avoids the associated dataset and inference cost. Relative to
RFID-based approaches~\cite{abhiram2025iotmonitor}, the proposed
system does not require per-product tagging or a cabinet-mounted
reader; per-product identification cost reduces to a printed
sticker. Relative to single-modality
prototypes~\cite{sarawathi2025spoilage,sarkar2024proactive,shetty2024monitoring},
the proposed system combines four modalities and exposes them
together through one mobile interface, which increases the
operator's awareness of cabinet state.

\subsection{Limitations}

\subsubsection{Lightweight-product sensitivity}
The load-cell subsystem is more reliable for items with a stable
bearing surface and a mass of approximately 50\,g or greater. Lighter
objects are sometimes missed or generate spurious events because
their signal falls within the noise band introduced by ambient
vibration and load-cell creep. A more rigid shelf, higher-grade load
cells, and a longer running-average filter would improve sensitivity
at the cost of some responsiveness.

\subsubsection{MQ135 baseline drift}
The MQ135 is useful as a trend indicator but does not provide
calibrated gas-concentration measurements. Its baseline drifts during
the first 30--60\,s after power-on. The firmware absorbs the
transient through an exponential moving average and the bridge
exposes a runtime endpoint that resets the baseline on demand. No
absolute concentration claims are made.

\subsubsection{Visual classifier scope}
The pixel-level classifier targets bananas only. Extending it to
other fruits would require revisiting the colour thresholds;
extending it to general food categories would benefit from a learned
model, which is a natural avenue for follow-on work. Sensitivity to
ambient lighting variation is partially mitigated by the LED-strip
illumination but remains a known constraint.

\subsubsection{Wi-Fi-only camera coupling}
The ESP32-CAM is reachable only when both the bridge host and the
camera share a local network. Deployments that cross network
boundaries would require either an SSH tunnel, a port forward, or a
Firebase-mediated relay. The Settings screen exposes the camera IP
explicitly so that the deployment can be reconfigured without
recompilation.

\subsubsection{Single-shelf prototype}
The reported prototype monitors one shelf. Multi-shelf operation
requires additional HX711 modules, each on its own \texttt{DT}/\texttt{SCK}
pin pair, or a chip-select arrangement. The firmware and bridge
changes required to scale to multiple shelves are minor but were
outside the scope of this demonstration.

\subsubsection{Firebase ships disabled}
The Flutter application carries a Firebase Realtime Database
integration that is wired in code but ships with placeholder
credentials. The deployed build operates in an in-memory mode by
default, which preserves reproducibility in environments without an
external Firebase project. A production deployment would replace the
placeholder \texttt{firebase\_options.dart} with credentials
obtained through \texttt{flutterfire configure}.

\subsubsection{Demonstration-scale evaluation}
The reported results were obtained from prototype-level
demonstration sessions rather than from controlled longitudinal
deployments. This work is positioned as a reproducible reference
design and a documented qualitative baseline rather than as a
finished evaluation; closing this gap is a natural direction for
future work.
